\documentclass[11pt]{article}
\usepackage[utf8]{inputenc}
\usepackage[T1]{fontenc}
\usepackage[italian]{babel}
\usepackage{geometry}
\geometry{a4paper, margin=1in}
\usepackage{xcolor}
\usepackage{listings}
\usepackage{hyperref}

\usepackage{tcolorbox}
\tcbuselibrary{breakable}
\begin{document}

\begin{tcolorbox}[colback=blue!5!white,colframe=blue!75!black,title=\textbf{DOMANDA}, breakable]
Temi d'esame (per ognuno ragionare sulla scelta del protocollo di livello trasporto)
\begin{itemize}
    \item Implementare un servizio di telepass: quando l'auto si avvicina al casello il terminale in esso presente spedisce al server il numero di targa; il server tiene traccia in una struttura dati del numero di passaggi per ogni targa e restituisce al client tale numero.
\end{itemize}
\end{tcolorbox}

\begin{tcolorbox}[colback=green!5!white,colframe=green!75!black,title=\textbf{RISPOSTA}, breakable]
\textbf{Piano di Implementazione:} Servizio Telepass (es19)

\textbf{1. Scelta del Protocollo di Trasporto}\\
La traccia richiede di inviare un numero di targa (stringa) al server, il quale deve incrementare un contatore per quella specifica targa e restituire il numero aggiornato di passaggi.

\textbf{Scelta raccomandata:} \texttt{TCP}

\begin{itemize}
    \item \textbf{Motivazione (Affidabilità e Modifica di Stato):} Il problema principale di questo scenario è che il server modifica uno stato interno (incrementa il numero di passaggi, che ha rilevanza per la fatturazione/tracciamento). Se usassimo \texttt{UDP} e il pacchetto di risposta (dal server al client) andasse perso, il client andrebbe in timeout e ritrasmetterebbe la richiesta. Il server, ricevendo nuovamente il datagramma \texttt{UDP}, incrementerebbe il contatore per la seconda volta per lo stesso passaggio fisico dell'auto, creando un'incongruenza grave (doppio addebito).
    \item \textbf{Vantaggio di \texttt{TCP}:} Usando \texttt{TCP}, la libreria di rete gestisce internamente le ritrasmissioni e l'eliminazione dei pacchetti duplicati. Il server è garantito di ricevere la richiesta una e una sola volta per ogni connessione stabilita dal casello per un dato veicolo. Se c'è un problema di rete, la connessione cade, ma non si creano addebiti "fantasma" dovuti a ritrasmissioni applicative non gestite.
    \item \textbf{Nota:} Si potrebbe usare \texttt{UDP} solo a patto di implementare a livello applicativo un identificatore univoco di transazione (ID Passaggio) per riconoscere i duplicati, ma in un contesto base \texttt{TCP} offre questa garanzia "gratis".
\end{itemize}
Pertanto, l'implementazione proposta utilizzerà \texttt{TCP}.

\textbf{2. Implementazione del Client (\texttt{clientTCP.c})}
\begin{itemize}
    \item Il client (Terminale Casello) instaura una connessione \texttt{TCP} con il server.
    \item Legge o simula l'acquisizione della targa dell'auto (es. \texttt{char targa[10]}).
    \item Invia la stringa della targa al server tramite \texttt{TCPSend}.
    \item Effettua una \texttt{TCPReceive} per attendere in risposta l'intero relativo al conteggio dei passaggi.
    \item Stampa il risultato (es. "L'auto AB123CD è passata 4 volte") e chiude la connessione.
\end{itemize}

\textbf{3. Implementazione del Server (\texttt{serverTCP.c})}
\begin{itemize}
    \item Crea un server \texttt{TCP} in ascolto.
    \item Gestisce una semplice struttura dati (es. un array di \texttt{struct} contenenti targa e conteggio, per semplicità limitato a un numero massimo di targhe).
    \item In un ciclo infinito, accetta una connessione da un casello (\texttt{acceptConnection}).
    \item Riceve la targa (\texttt{TCPReceive}).
    \item Cerca la targa nella struttura dati:
    \begin{itemize}
        \item Se esiste, incrementa il conteggio.
        \item Se non esiste, la aggiunge con conteggio = 1.
    \end{itemize}
    \item Risponde al client inviando il conteggio aggiornato (\texttt{TCPSend}).
    \item Chiude il socket del client e torna in attesa del prossimo passaggio.
\end{itemize}
\end{tcolorbox}

\end{document}
